% Source pp. 56-60 (Lectures7-12.pdf).
\section{Path Integral in Klein-Gordon Theory}

Let us now apply the idea of the (Euclidean) path integral to the KG field
theory. In this case the role of the generalized coordinates $q$ is taken by
the field degrees of freedom $\varphi(\xx)$, thus we have to replace
\[
q\to\varphi(\xx)\,,\qquad q(\tau)\to\varphi(\xx,\tau)=\varphi(x_E)\,,
\]
where $\tau$ is the ``imaginary time'' which I will denote $x_4$ below, and
$x_E=(\xx,x_4)$. The Euclidean action is obtained by substituting $t=-ix_4$ in
the KG action, i.e.
\[
\mathcal A=\int dx_4\,d\xx\,\frac12\Bigl[(\partial_4\varphi)^2+(\nabla\varphi)^2+m^2\varphi^2\Bigr]
=\int d^4x_E\,\left[\frac12\,(\partial_\mu\varphi)^2+\frac{m^2}{2}\,\varphi^2\right]\,,
\]
where now $\mu=1,2,3,4$. We are interested in the path integral
\[
Z=\int\bigl[D\varphi(x)\bigr]\,e^{-\mathcal A[\varphi(x)]}\,.
\]
More generally, we will be interested in the correlation functions defined as
the ratios
\begin{equation}\label{eq:corrdef}
\bigl\langle\varphi(x_1)\varphi(x_2)\cdots\varphi(x_N)\bigr\rangle
=\frac1Z\int\bigl[D\varphi\bigr]\,\varphi(x_1)\varphi(x_2)\cdots\varphi(x_N)\,e^{-\mathcal A[\varphi]}\,.
\end{equation}
Here and below I drop the subscript $E$ and use the notation $x_i$ for the
\emph{Euclidean} vectors, unless stated otherwise. As we discussed before, when
continued to complex values of $(x_i)_4$ and specialized to $(x_i)_4=it_i$,
these correlation functions yield the time-ordered expectation values
\[
\bra0\Tprod\bigl(\hat\varphi(\xx_1,t_1)\cdots\hat\varphi(\xx_N,t_N)\bigr)\ket0\,,
\]
where $\hat\varphi(\xx,t)$ denotes the real-time Heisenberg field operator.

Few words about the definition of the path integral in the case of fields. For
the path integral in quantum mechanics we had a definition in terms of a limit
$\Delta\tau\to0$ of a finite-dimensional integral with discretized ``time''
$\tau$. In principle, there are many ways to build finite approximations in
quantum field theory. For example, the 4-dimensional space can be replaced by,
say, the hypercubic lattice, in which $x^\mu$ take discrete values
\[
x^\mu=\Delta\,n^\mu\,,\qquad n^\mu\in\mathbb Z\,.
\]
The derivatives and the integration in the action can be replaced by finite
differences and summations, i.e.
\[
\partial_\mu\varphi(x)\to\frac1\Delta\,\bigl(\varphi(x+e^\mu\Delta)-\varphi(x)\bigr)\,;\qquad
\int d^4x\to\Delta^4\sum_x
\]
($e^\mu$ is the unit vector in the direction $\mu$); with this one can set
\begin{equation}\label{eq:latticemeasure}
\bigl[D\varphi\bigr]=\lim_{\Delta\to0}\prod_xd\varphi(x)\,.
\end{equation}
In the Klein-Gordon theory such definition works more or less
straightforwardly. In more interesting cases with an interaction defining the
integral over $[D\varphi]$ can become highly nontrivial problem, which is
closely related to the problem of \emph{critical behavior}; many subsequent
lectures will be devoted to this subject. Let me also note here that in the
context of field theory the integrals like those above are usually referred to
as \emph{functional integrals} (the term ``path integral'' is used
interchangeably).

The action $\mathcal A$ of the KG theory is quadratic in $\varphi$. This makes it
possible to evaluate the above path integrals explicitly. I am not going to
explain how these integrals are actually evaluated; this calculation is
presented in \textbf{PS, \S\,9.2}. Instead I will derive the above correlation
functions using formal properties of the functional integral. The resulting
formula is known as the

\subsection*{Wick's Theorem}

\begin{itemize}
\item The correlation function of any odd number of fields vanishes,
\[
\bigl\langle\varphi(x_1)\cdots\varphi(x_{2N+1})\bigr\rangle=0\,.
\]
This follows from the fact that the action $\mathcal A$ is an even functional
of $\varphi$, while in this case the product of an odd number of factors
$\varphi$ in the integrand is odd.
\item The two-point correlation function is
\[
\bigl\langle\varphi(x_1)\varphi(x_2)\bigr\rangle=D_E(x_1-x_2)\,,
\]
where $D_E(x)$ is the familiar function
\[
D_E(x)=\int\frac{d^4p}{(2\pi)^4}\,\frac1{p^2+m^2}\,e^{ipx}\,.
\]
From now on I will usually drop the subscript $E$, so that $D(x)$ will stand
for this function. Also, $x$ and $p$ now stand for the Euclidean vectors
$x^\mu=(\xx,x_4)$, $p^\mu=(\pp,p_4)$, with $xp:=\xx\cdot\pp+x_4p_4$.
\item The $2N$-point correlation function
\[
\bigl\langle\varphi(x_1)\cdots\varphi(x_{2N})\bigr\rangle
\]
is a sum of $(2N-1)!!\equiv1\cdot3\cdot\ldots\cdot(2N-1)$ terms, each term
corresponding to one of $(2N-1)!!$ distinct ways of ``pairing'' (or
``contracting'') between $2N$ field insertions
$\varphi(x_1)\cdots\varphi(x_{2N})$. A pairing of the insertions
$\varphi(x_i)$ with $\varphi(x_j)$ is usually denoted by an overbrace
\[
\dots\contraction{}{\varphi(x_i)}{\cdots}{\varphi(x_j)}\varphi(x_i)\cdots\varphi(x_j)\dots\,.
\]
Every contraction gives rise to the factor $D(x-x')$,
\[
\contraction{}{\varphi(x)}{\,}{\varphi(x')}\varphi(x)\,\varphi(x')=D(x-x')\,.
\]
\end{itemize}
The last rule of pairing can be expressed in a compact form, in terms of the
following

\subsection*{(Wick's) Recursion Relation}
\begin{equation}\label{eq:WickRec}
\bigl\langle\varphi(x_1)\varphi(x_2)\cdots\varphi(x_N)\bigr\rangle
=\sum_{k=2}^ND(x_1-x_k)\,\bigl\langle\varphi(x_2)\cdots\widehat{\varphi(x_k)}\cdots\varphi(x_N)\bigr\rangle\,,
\end{equation}
where the hat indicates that this insertion is dropped. This relation
expresses $N$-point correlation function in terms of the $N-2$-point ones.
When supplemented with the condition $\langle\varphi(x)\rangle=0$, it contains
all the statements of the Wick's theorem ($\langle1\rangle=1$ by definition).

Let us derive this relation. The main assumption will be the invariance of the
functional measure $D[\varphi]=D[\varphi']$ under the functional shifts
\[
\varphi(x)\to\varphi'(x):=\varphi(x)+c(x)
\]
with arbitrary fixed function $c(x)$. This invariance is an evident property
of the discrete approximation \eqref{eq:latticemeasure}, so we assume it.

Consider the functional integral representing the $N-1$-point correlation
function,
\[
Z^{-1}\int\bigl[D\varphi\bigr]\,\varphi(x_2)\cdots\varphi(x_N)\,e^{-\mathcal A[\varphi]}\,.
\]
Note that the product of the insertions starts with $\varphi(x_2)$, and
involves $N-1$ terms. Make the following change of the integration variables
\[
\varphi(x)\to\varphi'(x)=\varphi(x)+\epsilon(x)\,,
\]
where $\epsilon(x)$ is an arbitrary infinitesimal function. The action changes
as follows
\[
\mathcal A[\varphi+\epsilon]=\mathcal A[\varphi]+\int d^4x_1\,\epsilon(x_1)\,\bigl[m^2-\Delta_{x_1}\bigr]\varphi(x_1)+O(\epsilon^2)\,,
\]
where I have assumed that the function $\epsilon(x)$ decays sufficiently fast
making the integration by parts possible. Here
\[
\Delta_x:=\sum_{\mu=1}^4\left(\frac{\partial}{\partial x^\mu}\right)^2
\]
is 4D Laplace operator. Since the change of integration variables does not
change the value of the integral, we have the identity
\begin{multline*}
0=\sum_{k=2}^N\epsilon(x_k)\int\bigl[D\varphi\bigr]\,\varphi(x_2)\cdots\widehat{\varphi(x_k)}\cdots\varphi(x_N)\,e^{-\mathcal A}\\
-\int d^4x_1\,\epsilon(x_1)\,\bigl(m^2-\Delta_{x_1}\bigr)\int\bigl[D\varphi\bigr]\,\varphi(x_1)\varphi(x_2)\cdots\varphi(x_N)\,e^{-\mathcal A}\,.
\end{multline*}
As the function $\epsilon(x)$ is arbitrary, we must have
\begin{equation}\label{eq:WickSD}
\bigl(m^2-\Delta_{x_1}\bigr)\bigl\langle\varphi(x_1)\varphi(x_2)\cdots\varphi(x_N)\bigr\rangle
=\sum_{k=2}^N\delta^{(4)}(x_1-x_k)\,\bigl\langle\varphi(x_2)\cdots\widehat{\varphi(x_k)}\cdots\varphi(x_N)\bigr\rangle\,.
\end{equation}
Assuming that the $N-2$-point correlation functions in the right hand side are
known, we can think of this identity as of the equation for the $N$-point
correlation function in the left hand side. The solution has the form of our
recursion relation \eqref{eq:WickRec}, where the function $D(x-x')$ solves the
equation
\[
\bigl(m^2-\Delta_x\bigr)\,D(x-x')=\delta^{(4)}(x-x')\,.
\]

The Wick's theorem can be derived in many other ways. Direct evaluation of the
Gaussian functional integral is performed in the \textbf{Section 9.2 of PS}.
Derivation of equivalent rules in the operator formalism is explained in
\textbf{Sect.~4.3 of PS}.

The rules of the Wick's theorem are conveniently represented in terms of
diagrams. With every factor $D(x-x')$ one associates a line connecting the
points $x$ and $x'$,
\[
D(x-x')=\contraction{}{\varphi(x)}{}{\varphi(x')}\varphi(x)\varphi(x')
=\ x\;\dia{\draw[prop] (0,0) node[vtx]{} -- (3,0) node[vtx]{};}\;x'\,.
\]
The diagram consists of the lines representing the contractions. For example
\begin{align*}
\bigl\langle\varphi(x_1)\varphi(x_2)\varphi(x_3)\varphi(x_4)\bigr\rangle
&=\contraction{}{\varphi(x_1)}{}{\varphi(x_2)}\varphi(x_1)\varphi(x_2)\,
\contraction{}{\varphi(x_3)}{}{\varphi(x_4)}\varphi(x_3)\varphi(x_4)
+\ \text{two other ways of contractions}\\[4pt]
&=\figfile{wick4}
\end{align*}
Likewise, $2N$-point correlation function is represented as the sum of
$(2N-1)!!$ diagrams involving $N$ lines.
